\subsection{Frozen scientific question}
This preprint answers exactly one question (Q1):
\begin{quote}
`Is structural diagnosability a formally well-defined and decidable property of
enterprises modeled as continua (\KEA), independent of managerial intent and KPI
observability?''
\end{quote}

\subsection{Hard scope constraints}
We enforce the following constraints:
\begin{itemize}
  \item No new primitives beyond \OC{} core and ETS.K\_EA.v1.1.
  \item No reinterpretation of \OC{} or ETS.K\_EA.v1.1.
  \item No methodological or prescriptive content (no roadmaps, no consulting, no advice).
  \item No weakening of STOP / No-Go semantics and invariants.
  \item The paper stands alone; it does not require reading any whitepaper.
\end{itemize}

\subsection{Operational meaning}
The term `structural'' in this paper refers to structure encoded in the
enterprise continuum formalism (state space, boundaries, axes, thresholds,
flows, cycles, and continuity). The term `diagnosability'' refers to whether
a decision procedure can determine, from a formally admissible description of an
enterprise-as-continuum, whether the continuum is diagnosable in the sense
defined in \cref{sec:defns}.

We explicitly do \emph{not} assume the availability of KPIs, managerial intent,
or any specific human interpretation layer. Any observation model, if used, is
introduced only as a derived construct definable from the immutable primitives.
